\chapter*{Legacy Problem Recovery Audit for VI/05}
\addcontentsline{toc}{chapter}{Legacy Problem Recovery Audit for VI/05}

This editorial appendix records the disposition of the substantive legacy affine-geometry
problems relevant to irreducibility, components and connectedness.  The governing rule is
that a legacy problem is retained here, merged into an equivalent canonical dossier, or
explicitly routed to its mathematically correct chapter; it is never silently dropped.

\begin{longtable}{p{0.16\textwidth}p{0.34\textwidth}p{0.35\textwidth}}
\toprule
\textbf{Legacy item} & \textbf{Topic} & \textbf{Canonical disposition}\\
\midrule
DG4 P.1 & Zariski versus product topology & VI/02.P1; retained there\\
DG4 P.2 & Nonempty opens in irreducible sets & VI/05.P1; retained\\
DG4 P.3 & Distinguished opens form a basis & VI/02.P2; retained there\\
DG4 P.4 & Uniqueness of irreducible decomposition & VI/05.P2; retained\\
DG4 P.5 & Density of invertible/full-rank matrices & VI/02.P3; retained there\\
DG4 P.6 & Plane curves without common components; quasi-compactness & VI/05.P4; retained\\
DG4 P.7 & Hyperbola \(xy=1\), morphisms from \(\mathbb A^1\) & VI/04.P1; retained there\\
DG4 P.8 & Components of \(V(x^2-yz,xz-x)\) & VI/05.P3; retained\\
DG4 P.9 & Products of affine algebraic sets & VI/04.P2; retained there\\
DG4 P.10 & No nonconstant maps \(\mathbb A^1\to E\) & VI/04.P3; retained there\\
DG4 P.11 & Monomial curve \((t^3,t^4,t^5)\) & VI/03.P1; retained there\\
DG4 P.12 & Fibers and degeneration & VI/25; reserved for fibers/geometric fibers\\
DG4 P.13 & Unions of lines distinguished by local tangent data & VI/32; reserved for tangent spaces/local geometry\\
AG1 P.3 & Disconnectedness, idempotents, orthogonal idempotents, product rings & VI/05.P5; classical affine form retained; spectrum form returns later\\
AG5 P.3 & Integral iff reduced and irreducible & VI/27; scheme-level problem, not duplicated here\\
AG5 T05 selector & Scheme-level explanatory component/connectedness material & Used as a theory cross-check; no additional unique classical numbered problem was found in that selector\\
\bottomrule
\end{longtable}

\bigskip
\noindent\textbf{New canonical material.}
VI/05.P6--P8 and VI/05.C1--C3 are new synthesis items.  They do not replace legacy
problems; they deepen the component/minimal-prime/idempotent/graph dictionaries established
by the recovered corpus.

\bigskip
\noindent\textbf{Pedagogical rule.}
Examples are fully worked illustrations.  Exercises are shorter training items with
separate staged hints and full solutions.  Problems preserve harder synthesis and legacy
mathematics in self-contained dossiers.  Challenges are optional extensions.
